% Generated draft: requires mathematical and source review.
\documentclass[11pt]{article}
\usepackage[margin=25mm]{geometry}
\usepackage{fontspec}
\setmainfont{Latin Modern Roman}
\usepackage{amsmath,amssymb,mathtools,unicode-math}
\setmathfont{Latin Modern Math}
\usepackage{enumitem,needspace,xurl,hyperref}
\setlength{\parindent}{0pt}
\setlength{\parskip}{5pt}
\setlength{\emergencystretch}{3em}
\newcommand{\N}{\mathbb{N}}
\newcommand{\Z}{\mathbb{Z}}
\newcommand{\Q}{\mathbb{Q}}
\newcommand{\R}{\mathbb{R}}
\newcommand{\C}{\mathbb{C}}
\newcommand{\F}{\mathbb{F}}
\newcommand{\A}{\mathbb{A}}
\newcommand{\PP}{\mathbb{P}}
\newcommand{\E}{\mathbb{E}}
\newcommand{\eps}{\varepsilon}
\newcommand{\dd}{\,\mathrm{d}}
\begin{document}
\section*{UnsolvedMath ID: 30005387}
\textbf{Source problem code:} catalogue entry
\subsection{Definitions and mathematical statement}
A conventional interpretation uses oriented knots in the oriented sphere $S^3$. Fix a concordance category $\mathcal D$: either the smooth category or the locally flat topological category. Two knots are concordant in $\mathcal D$ if they cobound an annulus of that category in $S^3\times[0,1]$. Their concordance classes form the group $\mathcal C_{\mathcal D}$ under connected sum. A knot $K$ is slice in $\mathcal D$ if it bounds a disk of that category in $B^4$, equivalently $[K]=0$ in $\mathcal C_{\mathcal D}$. A class $[J]$ is torsion if
\[
\exists n\in\Z_{\geq 1}\quad n[J]=0.
\]
Thus torsion includes the zero class.

For a tubular neighborhood $\nu K$, let $\mu$ be its meridian and $\lambda$ its preferred longitude. Under the usual cabling convention, $K_{p,q}$ is the knot represented on $\partial\nu K$ by a primitive simple closed curve of class $p\lambda+q\mu$, so that $p$ is the longitudinal winding number and $\gcd(p,q)=1$. This specifies a conventional meaning of the notation; the supplied excerpt does not explicitly specify the slope convention.

\textbf{Supplied conjectural target.} With the category and cabling convention understood, the proposed assertion is that, for every knot $K$ and every pair of relatively prime integers $p,q$,
\[
[K_{p,q}]\text{ is torsion in }\mathcal C_{\mathcal D}
\quad\Longleftrightarrow\quad
\bigl(K\text{ is slice in }\mathcal D\ \text{and}\ |q|=1\bigr).
\]
Equivalently, the left side asks whether some positive number of connected-sum copies of $K_{p,q}$ is slice.

\textbf{Unresolved qualifications.} This display records the supplied target, rather than an unqualified, fully specified conjecture. The record does not identify the concordance category or explicitly state the ambient manifold. The conventional $S^3$ setting above therefore requires confirmation, and no assertion is made that the target is intended in both categories. The original excerpt allows relatively prime integers without restricting $p$ or excluding zero parameters. In particular, winding numbers $p=\pm1$ and the coprime zero-parameter pairs are not excluded. Under the convention above, $K_{1,1}$ is isotopic to $K$, so inclusion of this pair would make the proposed equivalence assert that every torsion concordance class is zero. The intended winding-number restriction must be checked; none has been silently inserted here.

The supplied excerpt states that $K$ slice and $q=1$ imply that $K_{p,q}$ is slice, attributes the converse question to Miyazaki, and introduces a stronger conjecture. The displayed torsion equivalence, including $|q|=1$, comes from the supplied clean statement; the original excerpt does not reproduce that stronger conjecture explicitly.
\subsection{Short English statement}
The proposed conjecture says that a cable knot has finite order in the knot concordance group exactly when its companion is slice and its meridional parameter has absolute value one. The supplied record leaves the concordance category and the intended range of cabling parameters unresolved, especially winding number one. These qualifications require human mathematical review.
\subsection{Sources}
\begin{itemize}
\item Supplied local problem registry, \texttt{data/unsolvedmath/problems.json}, permanent UnsolvedMath ID \(30005387\), source problem code \texttt{catalogue entry}, title \emph{Cable Torsion and Sliceness of Knots}. This draft uses the supplied complete record; no registry or external source was independently inspected.
\item \emph{Morphisms in Low Dimensions} (2023), Oberwolfach Reports, report 4/2023, pp.\ 215--259, DOI \texttt{10.4171/owr/2023/4}. Supplied source report and literature link: \url{https://doi.org/10.4171/owr/2023/4}. The record describes these reports as documenting mathematical workshops and their open-problem discussions.
\item The supplied excerpt identifies the contribution \emph{Involutive knot Floer homology and the Arf invariant} by Sungkyung Kang, joint work with JungHwan Park. It attributes the converse sliceness question to Miyazaki; no separate citation for that attribution is supplied. The registry has no named proposer, so authorship of the stronger conjecture is not inferred.
\item Kang and Park, \emph{Torsion in the knot concordance group and cabling}: \url{https://arxiv.org/abs/2207.11870}. This is a supplied literature reference, not an independently inspected source.
\item Literature assessment dated exactly \texttt{2026-08-22}: supplied dated curation, not independently verified. The record labels the problem \texttt{partially\_solved} and reports that the cable-torsion conjecture remains unsettled. It describes Kang and Park as constructing a nontrivial homomorphism detecting ``half'' of the torsion knots and obtaining infinite order for broad families of odd,\(1\) cables. The record supplies no precise hypotheses or interpretation of these summary phrases. Its background expressly calls this dated literature triage derived from the supplied Oberwolfach curation and requests independent verification of time-sensitive status and cited claims.
\item The registry labels the clean statement \texttt{corrected\_verified} and says it was checked against the official report for hypotheses, notation, and target. This is supplied source-verification metadata, not verification performed for this draft.
\end{itemize}
\end{document}
